\section{Machine Learning Co-Design}
\label{sec:ml-routing}

The low-latency guarantees of \textsc{Nexus} rely on a hybrid routing system that avoids invoking the heavy generative LLM for high-confidence decisions. This is achieved through a machine learning co-design consisting of representation engineering, statistical margin calibration, and adversarial cross-encoder training.

\subsection{Representation Engineering: Transient Intent Signatures}
Dense retrieval models (e.g., Nomic-Embed) project text documents into a latent vector space $\mathbb{R}^D$. In tool routing, user queries are short and colloquial, whereas tool schemas are long, structured, and formal. This syntactic discrepancy causes low similarity scores and poor separation in the embedding space.

To resolve this, \textsc{Nexus} implements \textbf{Transient Intent Signatures}. Rather than prepending intent descriptors to user queries online (which could introduce query latency and lexical distortion), intent signatures (verb phrases describing the tool's action) are prepended to the \textbf{tool schema document} $d_{\text{tool}}$ offline prior to embedding generation. The online user query merely receives a static Nomic retrieval prefix (\texttt{"search\_query: "}). In the latent embedding space, this prepended signature induces a translation vector $\vec{t}_{\text{intent}}$, displacing the tool description representation along intent-specific manifolds:
\begin{equation}
E_{\text{engineered}}(s_{\text{intent}} \oplus d_{\text{tool}}) \approx E(d_{\text{tool}}) + \vec{t}_{\text{intent}}
\end{equation}
where $\oplus$ denotes string concatenation. We emphasize that this formulation represents an empirical representational heuristic in the latent space of deep networks rather than a strict algebraic identity. This displacement separates embedding representations of distinct tools in the latent dense space, increasing the distance between the target tool cluster and adjacent non-target clusters, thereby improving dense retrieval recall.

\subsection{Statistical Margin Calibration and the $P_{20}$ Gate}
While dense retrieval (SLB) is computationally efficient, it exhibits a tail of misclassifications under paraphrased or adversarial queries. A heavy Cross-Encoder (CE) model can rerank candidate tools to resolve ambiguity, but invoking it on every request introduces millisecond-scale overhead, violating the routing latency budget.

To bound the computational cost, we implement a statistical decision gate based on the confidence margin. Let $\text{Score}_{\text{Top1}}$ and $\text{Score}_{\text{Top2}}$ be the cosine similarity scores of the highest-scoring tool candidates retrieved by the L1 SLB. The confidence margin $\Delta$ is defined as:
\begin{equation}
\Delta = \text{Score}_{\text{Top1}} - \text{Score}_{\text{Top2}}.
\end{equation}

If $\Delta \ge \tau$, the classification is deemed high-confidence, and the MMU executes the Auto-Routing Fast Path, bypassing the Cross-Encoder. If $\Delta < \tau$, the classification is ambiguous, and the Cross-Encoder is invoked to rerank the candidates.

To extract the threshold $\tau$, we evaluate the empirical cumulative distribution function (CDF) of the margin $\Delta$ over a validation set of adversarial queries:
\begin{equation}
F_{\Delta}(x) = P(\Delta \le x).
\end{equation}
We statically calibrate the threshold at the $20\text{th}$ percentile ($P_{20}$) of the margin distribution:
\begin{equation}
\tau = F_{\Delta}^{-1}(0.20) \approx 0.01365.
\end{equation}
By setting $\tau \approx 0.01365$, we strictly bound the Cross-Encoder invocation rate to a maximum $20\%$ budget under production workloads, ensuring that $80\%$ of requests execute via the microsecond-scale auto-routing fast path.

\subsection{Synthetic Adversarial Training of CE v3}
To prevent data leakage and guarantee robustness under diverse user phrasing, the Cross-Encoder reranker (CE v3) is trained strictly on synthetic adversarial datasets. We generate a corpus of paraphrased queries using an offline generator, applying semantic distortions, grammatical variations, and spelling perturbations. 

The training objective enforces a margin-ranking loss over positive and negative query-tool pairs, utilizing a hard-negative mining strategy to separate highly similar function schemas. By isolating the training corpus from the evaluation hold-out sets, we prevent overfitting and guarantee robust out-of-distribution generalization in multi-domain registries.
